Prompt tone is the most freely-varying part of any LLM interaction, yet there is little quantitative evidence on whether \emph{skeptical} tone (separate from skeptical \emph{content}) helps or hurts reasoning. Prior work treats user skepticism as monotonically harmful — flipping correct answers and triggering apologetic capitulation — while the prompt-engineering literature shows that emotional and motivational stimuli can boost accuracy. We test whether a graded preemptive judgmental tone applied \emph{before} the question opens a ``sweet spot'' between these two regimes. Across 3{,}900 single-turn and 1{,}200 follow-up API calls covering two frontier OpenAI models, three datasets (\tqa, \gsm, \mmlu), and a content-fixed six-level tone ladder (neutral $\rightarrow$ hostile), we find three results. First, models think \emph{longer} when judged: completion tokens grow by 19--57\% from neutral to strongly-judgmental on \tqa, significant at $p \le 10^{-9}$ (Wilcoxon, paired). Second, a small directional accuracy lift appears on the one task with headroom (\tqa, $+3$--$4$\,pp) but is not significant at our per-cell sample sizes; \gsm and \mmlu sit at ceiling. Third, the explicit apology marker the sycophancy literature uses for capitulation detection is \emph{exactly zero} across all $5{,}100$ responses — yet preemptive judgmental tone makes the model \emph{more} fragile to a follow-up rebuttal: post-rebuttal correct-to-wrong flips on \tqa roughly triple from $4\%$ at neutral to $13\%$ at judgmental. The single-turn sweet spot is real but not free; multi-turn settings invert it.
